\section*{الفصل \ref{ch:b2:guided-waves} --- الأمواج الموجَّهة والأجواف}

\begin{solution}{exo:b2:guided-waves:1}
$f_n = nc/2a = 7.5$، $15$، \qty{22.5}{GHz}. وعند \qty{12}{GHz} ينتشر $n = 1$ وحده: $k_g =
(2\pi/c)\sqrt{f^2 - f_1^2} = \qty{196}{rad/m}$، $\lambda_g = \qty{3.2}{cm}$، $v_\varphi = c/0.78
= \qty{3.8e8}{m/s}$، $v_g = \qty{2.3e8}{m/s}$. وعند \qty{6}{GHz}: $\ell = c/2\pi\sqrt{f_1^2 -
f^2} = \qty{1.1}{cm}$.
\end{solution}

\begin{solution}{exo:b2:guided-waves:2}
TE$_{10}$ عند \qty{6.55}{GHz}، و TE$_{20}$ عند \qty{13.1}{GHz}، و TE$_{01}$
عند \qty{14.7}{GHz}: فالنمط واحد من $6.55$ إلى \qty{13.1}{GHz}. وعند \qty{10}{GHz}: $\sqrt{1 - (6.55/10)^2} =
0.756$، $\lambda_g = \qty{4.0}{cm}$، $v_g = 0.76c$، $\cos\theta = 0.756$، $\theta = \qty{41}{\degree}$.
و $b \approx a/2$ يبقي TE$_{01}$ فوق TE$_{20}$ (أي أعرض نطاق أحادي النمط)
مع ترك أوسع فرجة للحقل (أي الاستطاعة قبل الانهيار).
\end{solution}

\begin{solution}{exo:b2:guided-waves:3}
$f = 1.5 \times 10^8\sqrt{(m/0.3)^2 + (n/0.3)^2 + (p/0.2)^2}$: \qty{0.71}{GHz}، و \qty{0.90}{GHz}،
و \qty{1.35}{GHz}، و \qty{1.95}{GHz}. و $N = 8\pi Vf^3/3c^3 = 8\pi \times 0.018 \times 1.56 \times 10^{28}/
8.1 \times 10^{25} \approx 260$. وأما المكعب: فيكون $f = c\sqrt2/2L$، و $L = \qty{2.3}{cm}$.
\end{solution}

\begin{solution}{exo:b2:guided-waves:4}
$\mathrm{NA} = \sqrt{1.48^2 - 1.46^2} = 0.24$، \qty{14}{\degree}؛ و $\Delta t = (n_1L/c)(n_1/n_2
- 1) = \qty{68}{ns/km}$؛ و \qty{135}{ns} على \qty{2}{km}: \qty{7}{Mbit/s}؛ ومع $0.003$:
\qty{10}{ns/km}، \qty{50}{Mbit/s}.
\end{solution}

\begin{solution}{exo:b2:guided-waves:5}
(أ) $\underline B_y = (k_g/\omega)\underline E$، و $\underline B_z = (\iu/\omega)\partial_y\underline E = (\iu\pi/a\omega)
E_0\cos(\pi y/a)\eu^{\iu(\omega t - k_gz)}$: أي في التربيع. (ب) وعند $y = 0$ يكون $B_y = 0$ و
$\vect j_s = \vect e_y\wedge\vect B/\mu_0 = (B_z/\mu_0)\vect e_x$: أي وفق الحقل. (ج)
و $\langle\Pi_z\rangle = E_0^2k_g\sin^2(\pi y/a)/2\mu_0\omega$، والاستطاعة في وحدة العرض
$E_0^2k_ga/4\mu_0\omega$؛ و $\langle\Pi_y\rangle \propto \operatorname{Re}(\underline E\,\iu\underline B_z^*)
= 0$. (د) ومتوسط الطاقة في وحدة الطول والعرض $\tfrac14\varepsilon_0E_0^2a$ (والنصفان
الكهربائي والمغناطيسي كلٌّ منهما $\tfrac18\varepsilon_0E_0^2a$)؛ والنسبة هي $c^2k_g/\omega = v_g$.
\end{solution}

\begin{solution}{exo:b2:guided-waves:6}
(أ) $f_1 = \qty{150}{GHz} \gg f$: $\ell \approx a/\pi = \qty{0.32}{mm}$؛ و $\eu^{-0.5/0.32}$:
\qty{-14}{dB} في السعة، و $-27$ في الاستطاعة. (ب) $f_1 = \qty{750}{MHz}$؛ وعند \qty{100}{MHz}
يكون $\ell = \qty{6.4}{cm}$: \qty{-136}{dB} على متر. (ج) $f_1 = \qty{15}{GHz}$، و $\ell =
\qty{3.2}{mm}$: \qty{27}{dB/cm}. (د) لأن الموجة لا تسع؛ فالتلاشي
هندسة لا تبديد.
\end{solution}

\begin{solution}{exo:b2:guided-waves:7}
(أ) $P = \tfrac14\varepsilon_0cE_0^2ab\sqrt{1 - f_{10}^2/f^2} = \qty{120}{kW}$. (ب) و $\times9$:
\qty{1}{MW}؛ والآزوت المضغوط أو SF$_6$ يرفعان حقل الانهيار. (ج)
و $j_s \approx E_0/\mu_0c = \qty{2.7}{kA/m}$؛ و $\tfrac12R_sj_s^2 = \qty{9e4}{W/m^2}$ على $2a =
\qty{0.046}{m^2}$ في المتر: \qty{4}{kW/m}، أي $3.5\%$ في المتر، ونحو \qty{0.15}{dB/m}.
(د) أي أقلّ من الكابل بثلاث مرات، ومع ميغاواطات بدل
كيلوواطات.
\end{solution}

\begin{solution}{exo:b2:guided-waves:8}
(أ) $f = c\sqrt2/2L = \qty{2.12}{GHz}$. (ب) $\delta = \qty{1.4}{\micro m}$، و $R_s = \qty{12}{m\ohm}$؛
و $W = \varepsilon_0E_0^2V/8 = \qty{1.1e-15}{}E_0^2$؛ و $P = \tfrac12R_s(E_0/377)^2 \times 6L^2 =
\qty{2.5e-9}{}E_0^2$: فيكون $Q = \omega W/P \approx 6000$. (ج) \qty{360}{kHz}؛ و \qty{0.45}{\micro s}. (د)
و $Q \sim 10^{10}$: فالحقل البالغ عشرات الميغافولطات في المتر تديمه
كيلوواطات بدل غيغاواطات.
\end{solution}

\begin{solution}{exo:b2:guided-waves:9}
(أ) النقاط $(m, n, p)$ التي $m^2 + n^2 + p^2 \le (2fL/c)^2$ تملأ ثُمن
كرة: $\tfrac18\cdot\tfrac43\pi(2fL/c)^3$، مضروبًا في استقطابين:
$8\pi L^3f^3/3c^3$. (ب) $8\pi f^2/c^3$. (ج) $9 \times 10^{-7}$ في الهرتز (أي واحد في
الميغاهرتز) عند \qty{1}{GHz}؛ و $2 \times 10^5$ في الهرتز عند \qty{500}{THz}. (د) والأنماط
منفصلة حين يتجاوز تباعدها عرضها؛ وفي الفرن نمط
كل \qty{10}{MHz} تقريبًا حول \qty{2.45}{GHz} --- أي عشرات ضمن نسبة
قليلة.
\end{solution}

\begin{solution}{exo:b2:guided-waves:10}
(أ) من $\operatorname{\vect{curl}}\vect E = -\partial_t\vect B$:
\[
\underline B_x = -\frac{k_g}\omega\underline E , \qquad
\underline B_z = \frac{\iu\pi}{a\omega}E_0\cos\frac{\pi x}a\,\eu^{\iu(\omega t - k_gz)} ,
\]
و $\partial_xB_x + \partial_zB_z = -(k_g\pi/a\omega)E_0\cos(\pi x/a) + (k_g\pi/a\omega)E_0
\cos(\pi x/a) = 0$. (ب) و $(\operatorname{\vect{curl}}\vect B)_y = \partial_zB_x - \partial_xB_z =
(\iu/\omega)(k_g^2 + \pi^2/a^2)\underline E = (\iu\omega/c^2)\underline E$: أي $k_g^2 = \omega^2/c^2 - \pi^2/a^2$.
(ج) الجداران العريضان: $\vect j_s = \pm(B_z\vect e_x - B_x\vect e_z)/\mu_0$ --- أي وفق $z$
على صورة $\sin(\pi x/a)$، وعرضًا على صورة $\cos(\pi x/a)$ (وهو معدوم في الوسط)؛ والجداران
الضيقان: وفق $y$ فقط. (د) والشقّ الطولي المركزي لا يقطع إلا
التيارات المستعرضة، وهي منعدمة هناك؛ أما الشقّ المستعرض فيقطع
الطولية ويشعّ: وهو هوائي \omterm{prop:b2:guided-waves:plates}{دليل الموجة} المشقوق.
\end{solution}

\begin{solution}{exo:b2:guided-waves:11}
(أ) $2a\sqrt{n_1^2 - n_2^2}/\lambda = 2 \times 50 \times 0.242/1.5 \approx 16$ في كل استقطاب.
(ب) $a < \lambda/2\mathrm{NA} = \qty{3.1}{\micro m}$. (ج) ولبٌّ دائري فتحته $\mathrm{NA} =
0.12$ أحادي النمط دون $d = 2.405\lambda/\pi\mathrm{NA} = \qty{10}{\micro m}$: أي
معيار \qty{9}{\micro m}. (د) والشعاع المماسّ ينعكس كليًّا دائمًا: فأدنى
نمط موجود عند كل تواتر، وذيوله المتلاشية لا تفعل أكثر من
الامتداد في الغلاف.
\end{solution}

\begin{solution}{exo:b2:guided-waves:12}
(أ) الشعاع الخارج عن المحور يقضي زمنه في قرينة أدنى، حيث الضوء
أسرع: فيُعوَّض المسار الأطول. (ب) و $n(r)\cos\theta = n_1\cos\theta_0$
مع $n \approx n_1(1 - \Delta r^2/a^2)$ و $\cos\theta \approx 1 - \theta^2/2$: $(\dd r/\dd z)^2 =
\theta_0^2 - 2\Delta r^2/a^2$، وهو تذبذب توافقي $r = r_0\sin(\sqrt{2\Delta}\,z/a)$.
(ج) $2\pi a/\sqrt{2\Delta} = \qty{1.1}{mm}$. (د) و $n_1L\Delta^2/2c = \qty{0.25}{ns/km}$ مقابل
\qty{49}{ns/km}: أي أحسن بمئتي مرة.
\end{solution}

\begin{solution}{pb:b2:guided-waves:1}
\textbf{1.} $c/2a = \qty{1.74}{GHz}$؛ و TE$_{20}$ و TE$_{01}$: \qty{3.49}{GHz}: فلا ينتشر إلا
TE$_{10}$ عند \qty{2.45}{GHz}.

\textbf{2.} $\sqrt{1 - (1.74/2.45)^2} = 0.70$: $\lambda_g = \qty{17.4}{cm}$، $v_\varphi = 1.42c$،
$v_g = 0.70c = \qty{2.1e8}{m/s}$.

\textbf{3.} $E_0^2 = 4P/\varepsilon_0c\,ab\sqrt{\cdot} = 3200/6.9 \times 10^{-6}$: $E_0 = \qty{21}{kV/m}$،
أي جزء من مئة من الانهيار.

\textbf{4.} $j_s \approx E_0/\mu_0c = \qty{57}{A/m}$؛ و $\delta = \qty{10}{\micro m}$، و $R_s = \qty{0.1}{\ohm}$؛
و $\tfrac12R_sj_s^2 \times 2a = \qty{28}{W/m}$: أي $3.5\%$ في المتر --- فالفولاذ ضيّاع،
والدليل قصير.

\textbf{5.} $f = 1.5 \times 10^8\sqrt{11.1(m^2 + n^2) + 25p^2}$ ضمن $\pm3\%$ من
\qty{2.45}{GHz}: $(1, 2, 3)$، $(2, 1, 3)$ عند \qty{2.51}{GHz}، و $(4, 0, 2)$، $(0, 4, 2)$،
$(4, 3, 0)$، $(3, 4, 0)$ عند \qty{2.50}{GHz} --- أي نحو ستة.

\textbf{6.} أنصاف أطوال الموجة $a/m$، $b/n$، $d/p$: أي \omterm{def:b2:maxwell-equations:operators}{تباعد} من $7$ إلى \qty{10}{cm}. والصينية
الدوّارة تجرّ الطعام عبر مواضع العظم والصغر؛ والمقلّب (وهو
مروحة معدنية دوّارة) يعيد خلط الأنماط.

\textbf{7.} $f/Q = \qty{12}{MHz}$: فتتراكب الرنينات المحمَّلة في
مستمرّ؛ والمغنيترون يجد دائمًا نمطًا يغذّيه.

\textbf{8.} $W = QP/\omega = 3000 \times 800/1.54 \times 10^{10} = \qty{0.16}{J}$؛ و $u = \qty{8.7}{J/m^3}$؛
و $E_0 \approx \sqrt{4u/\varepsilon_0} = \qty{2}{MV/m}$ --- أي قرب الانهيار: فأقواس؛ والاستطاعة
المنعكسة ترجع إلى المغنيترون، وهو محميّ (قليلًا) بحمل
صوري أو بمدوِّر --- ومن هنا "لا تشغّله فارغًا".

\textbf{9.} $0.2/2.1 \times 10^8 = \qty{1}{ns}$.

\textbf{10.} \qty{1.6}{kW} في الذروة؛ وبين الرشقات يخفت الحقل في
$Q/\omega = \qty{13}{ns}$: فالجوف فارغ في معظم الوقت.

\textbf{11.} $\ell \approx a/\pi = \qty{0.48}{mm}$؛ و $\eu^{-1/0.48}$: \qty{-18}{dB} في السعة،
و \qty{-36}{dB} في الاستطاعة.

\textbf{12.} $\lambda/4 = \qty{3.1}{cm}$.

\textbf{13.} \qty{5}{mW/cm^2} $\times$ \qty{100}{cm^2} = \qty{0.5}{W}: أي $0.06\%$.

\textbf{14.} $f_1 \approx c/2d = \qty{3.75}{GHz} > \qty{2.45}{GHz}$: أي دون القطع؛ و $\ell =
\qty{1.7}{cm}$، و $\eu^{-30/1.7}$: \qty{-150}{dB}.

\textbf{15.} يجب أن تجري تيارات الجدران متصلة إلى الشبكة؛
والفجوة في الوصل شقٌّ يقطعها --- أي هوائي شقّ
يشعّ إلى الخارج.

\textbf{16.} $\mathrm{NA} = \sqrt{1.465^2 - 1.460^2} = 0.12$، \qty{7}{\degree}: أي مخروط
ضيق، يحتاج إلى ليزر وعدسة.

\textbf{17.} $(n_1L/c)(n_1/n_2 - 1) = \qty{17}{ns/km}$؛ و \qty{170}{ns} على \qty{10}{km}:
\qty{6}{Mbit/s}.

\textbf{18.} $\tfrac12(\pi d\,\mathrm{NA}/\lambda)^2 = 75$؛ ومن أجل \qty{9}{\micro m}: $2.4$ --- أي
النظام الأحادي النمط ($V = \pi d\,\mathrm{NA}/\lambda = 2.2 < 2.4$).

\textbf{19.} $\Delta x_0 = (c/n)\tau = \qty{1}{cm}$: $t_{\text{تمدد}} = 10^{-4}/0.2 = \qty{0.5}{ms}$،
و \qty{100}{km}: أي \qty{10}{Gbit/s} على تلك المسافة قبل التعويض.

\textbf{20.} \qty{20}{dB}، أي عامل $100$: \qty{10}{\micro W}؛ و $6000/80 = 75$
مضخّمًا.

\textbf{21.} $f = \qty{1.9e14}{Hz}$، و $hf = \qty{0.8}{eV}$؛ و $8 \times 10^{15}$ فوتون في
الثانية؛ و $8 \times 10^5$ في البتّ عند الدخل، و $8 \times 10^3$ بعد \qty{100}{km}.

\textbf{22.} مجموع الضياعين أصغر ما يكون قرب \qty{1.55}{\micro m}
(\qty{0.2}{dB/km}).

\textbf{23.} $R = ((1.465 - 1)/2.465)^2 = 3.6\%$: أي \qty{0.16}{dB} في كل وجه، ومرتان
عند الوصلة زائدًا تداخل الفجوة؛ والهلام أو التماسّ الفيزيائي
يزيل الهواء.

\textbf{24.} عند انحناء لطيف يظل الشعاع يلاقي الحدّ بعد
الزاوية الحرجة؛ وعند انحناء حادّ يهبط الورود على الجدار الخارجي
دونها فيتسرب الضوء.

\textbf{25.} الفرن: جدران معدنية، وانعكاس $r = -1$؛ وعرض النطاق تحدده
أنماط الجوف؛ والضياعات في قشرة الفولاذ. والليف: \omterm{thm:b2:wave-interfaces:snell}{انعكاس كلي
داخلي} عند حدّ زجاج--زجاج؛ وعرض النطاق يحدده \omterm{prop:b2:guided-waves:fibre}{التبدد النمطي}
واللوني؛ والضياعات بالتبعثر والامتصاص، \qty{0.2}{dB/km}.
\end{solution}
